\section{Related Work}
\label{sec:related}

\textbf{Encoder--decoder summarization.} Abstractive summarization inherits its
architecture from neural machine translation. Sutskever et al.
\cite{sutskever2014seq2seq} showed that a recurrent encoder could compress a
variable-length input into a fixed-length vector from which a decoder generates
the output. That fixed-length vector is also the approach's weakness: a single
state must carry an entire document, and the longer the document the more it
must discard. Bahdanau et al. \cite{bahdanau2015attention} removed the
bottleneck by letting the decoder attend over all encoder states, learning a
soft alignment between output and input positions. Vaswani et al.
\cite{vaswani2017attention} then discarded recurrence entirely in favour of
self-attention. Within summarization specifically, See et al.
\cite{see2017pointer} addressed two failure modes of the recurrent formulation
--- inability to reproduce rare words, and repetition --- with a pointer
mechanism that copies from the source and a coverage penalty. Our recurrent
baseline follows this lineage, using additive attention
\cite{bahdanau2015attention}, though we did not train it (Section~\ref{sec:limitations}).

\textbf{Pretraining and its benchmarks.} Pretraining changed the cost structure
of the task: a general model is adapted with comparatively little in-domain
data. T5 \cite{raffel2020t5} casts every task as text-to-text and selects among
them with a textual prefix, which is the model and the mechanism we use. Liu
and Lapata \cite{liu2019bertsum} adapted a pretrained encoder to both
extractive and abstractive summarization, and PEGASUS \cite{zhang2020pegasus}
introduced a pretraining objective designed for summarization specifically.

These systems are, with few exceptions, developed and reported on news
corpora, where a reference summary is a paragraph or --- in the most compressed
case --- a single sentence. Our references are neither. They are headlines a
customer typed into a form, averaging four words. This is not a minor
difference of degree: it changes what the standard metric can measure, which we
take up in Section~\ref{sec:analysis}. We stress it here because it is the main
reason absolute ROUGE values in this paper should not be read against the
numbers those systems report.

\textbf{Opinion summarization.} Summarizing many reviews of one product is a
distinct problem, and it is usually attacked without reference summaries, since
no customer writes a summary of everyone else's opinions. MeanSum
\cite{chu2019meansum} trains an unsupervised model to generate a single review-
like summary whose encoding stays close to the mean of the input reviews.
Bra\v{z}inskas et al. \cite{brazinskas2020fewshot} showed that a small number of
annotated summaries is enough to adapt such a model. Both produce one fluent
paragraph per product. We differ in output rather than ambition: we combine
per-review summaries through a map--reduce procedure and pair them with a
structured list of praised and criticised aspects, because a fluent paragraph
is precisely the format in which a minority complaint disappears --- the effect
we document in Section~\ref{sec:analysis}.

\textbf{Evaluation.} ROUGE \cite{lin2004rouge} remains the default and was
designed against multi-sentence references. Its limits are well documented:
BERTScore \cite{zhang2020bertscore} replaces surface overlap with contextual
embedding similarity, and SummEval \cite{fabbri2021summeval} found that
automatic metrics, ROUGE included, correlate weakly with human judgement.
Maynez et al. \cite{maynez2020faithfulness} showed further that overlap metrics
are largely blind to faithfulness, since a summary can score well while
asserting something the source does not support. We report ROUGE, and
Section~\ref{sec:analysis} sets out what we believe it does and does not
establish in our setting.
